\chapter{课程体系规划与全球教学实践}

\section{三大能力领域的课程打包整合}

\begin{preamblebox}[从孤立课程走向能力集聚]
在 CS2023 指导纲要中，知识领域（KA）是内容切片的逻辑集合，而非孤立的高校课程。为了避免传统教学中“各自为政、知识脱节”的弊病，工作组围绕产业界与研究界的三大核心能力支柱，提出了跨领域的课程打包与贯穿整合范式：
\end{preamblebox}

\subsection{软件核心能力领域}
涵盖软件构建与理论分析的全生命周期，深度整合以下四大知识领域：
\begin{itemize}
    \item \textbf{软件开发基础 (SDF)}：编程思维、算法设计初步、防御性编码与调试。
    \item \textbf{算法与复杂性 (AL)}：高阶抽象数据类型、图算法、NPC 理论与算法伦理。
    \item \textbf{编程语言理论基础 (FPL)}：类型系统、形式语义、函数式/面向对象范式与运行时垃圾回收。
    \item \textbf{软件工程 (SE)}：大型代码仓演进、CI/CD 持续集成、重构、敏捷迭代与架构质量。
\end{itemize}
\textbf{【典型整合课程】}：《高级软件构造与算法设计》、《面向对象系统分析与形式化验证》、《企业级软件架构演化实训》。

\subsection{系统核心能力领域}
聚焦软硬件交界面及大规模分布式基础设施，深度整合以下七大知识领域：
\begin{itemize}
    \item \textbf{系统基础 (SF)}：系统思维、软硬件协同、状态机与资源抽象。
    \item \textbf{体系结构与组织 (AR)}：微架构、流水线、ILP、缓存一致性与 DSA/PIM 加速器。
    \item \textbf{操作系统 (OS)}：进程线程并发、虚拟内存、驱动架构与安全隔离。
    \item \textbf{并行与分布式计算 (PDC)}：OpenMP/MPI 并行、GPU 异构计算、分布式共识。
    \item \textbf{网络与通信 (NC)}：协议栈、TCP/QUIC 拥塞控制、BGP 路由与数据中心 Spine-Leaf 网络。
    \item \textbf{计算机安全 (SEC)}：内存防护、权限提升防御、威胁建模与系统韧性。
    \item \textbf{数据管理 (DM)}：底层存储引擎、WAL 预写日志、事务 ACID 与 LSM-Tree 存储结构。
\end{itemize}
\textbf{【典型整合课程】}：《现代计算机系统软硬件全栈设计》、《高并发分布式云原生系统》、《操作系统与网络协议栈内核实训》。

\subsection{应用核心能力领域}
面向真实应用场景，深度整合以下六大知识领域：
\begin{itemize}
    \item \textbf{人工智能 (AI)}：机器学习基础、深度神经网络、大语言模型（LLM）工程与强化学习。
    \item \textbf{图形学与交互技术 (GIT)}：现代实时渲染管线、物理仿真、XR/VR/AR 空间计算。
    \item \textbf{人机交互 (HCI)}：以人为本的设计（HCD）、可用性工程度量与多模态交互体验。
    \item \textbf{专用平台开发 (SPD)}：移动生态、物联网微控制器与边缘计算。
    \item \textbf{计算机安全 (SEC)} 与 \textbf{数据管理 (DM)}：应用层隐私注入、SQL 审计与多源数据管道。
\end{itemize}
\textbf{【典型整合课程】}：《智能具身感知与人机协同系统》、《空间计算与三维数字孪生引擎》、《数据驱动智能应用工程》。

\section{高校专业培养方案制定建议}

针对不同高校的定位与办学特长，CS2023 提出因校制宜的培养方案规划路径：
\begin{enumerate}
    \item \textbf{软件与工程主导型培养方案}：夯实编程底座（SDF）与算法深度（AL），在大三、大四重点强化软件工程工业级流水线，注重代码评审、自动化测试覆盖率与高内聚低耦合系统重构。
    \item \textbf{系统与底层架构主导型培养方案}：打通“数字逻辑 $\to$ 体系结构 $\to$ 操作系统 $\to$ 网络 $\to$ 分布式系统”全链路，强调动手编写简易 CPU、实现微内核与自研网络协议栈，培养具备底层破局能力的高精尖系统架构人才。
    \item \textbf{智能与交叉应用主导型培养方案}：紧扣现代 AI、大模型与沉浸式人机交互，前置连续数学基础（线性代数与概率优化），让学生在具体垂直行业（如智慧医疗、自动驾驶、金融科技）中开展实战课题。
\end{enumerate}

\section{全域高阶能力任务规范}

\begin{taskbox}[跨领域重大任务 1：高可用分布式智能交易结算平台架构设计]
\textbf{【跨领域融合】}：Systems + Software + Applications（涵盖 PDC, OS, NC, SEC, DM, SE）\\
\textbf{【能力陈述】}：针对千万级 TPS 实时交易结算场景，综合运用分库分表、分布式事务（Saga/两阶段提交）、Raft 共识日志强同步、零信任网络加密及全链路高可用监控，设计系统整体架构，制定故障自动容灾转移演练方案。\\
\textbf{【达标等级】}：\textbf{开发与评估}
\end{taskbox}

\begin{taskbox}[跨领域重大任务 2：面向低功耗边缘端的大模型轻量化推理引擎]
\textbf{【跨领域融合】}：Systems + Applications（涵盖 AR, AI, SPD, SF）\\
\textbf{【能力陈述】}：在算力与电量受限的嵌入式/移动硬件平台上部署大语言模型，通过权重量化（INT4/FP8）、算子融合、剪枝及硬件特定 NPU 汇编级指令调优，实现毫秒级首字延迟与超低功耗运行。\\
\textbf{【达标等级】}：\textbf{开发}
\end{taskbox}

\section{全球计算机教育前沿实践与时代思考}

\subsection{社会价值与伦理融合实践}
\begin{itemize}
    \item \textbf{无障碍辅助技术教学}：将 WCAG 标准、无障碍屏幕阅读器、色障友好交互作为前端与系统设计的核心质量指标。
    \item \textbf{科技向善}：引导学生参与非营利组织开源项目，将算法算力用于环境保护、贫困消减与公共卫生。
    \item \textbf{全球伦理视野与南北半球互通}：不仅探讨欧美发达经济体的技术法律框架，更深入审视发展中国家在数据主权、数字鸿沟与算法剥削上的严峻现实。
\end{itemize}

\subsection{全球化视野下的计算机教育模式}
\begin{itemize}
    \item \textbf{中国计算机教育实践}：探讨中国“101 计划”、新工科建设、拔尖创新人才培养基地以及产教融合协同育人机制，为全球超大规模本科计算机教学质量提升提供中国方案。
    \item \textbf{通识博雅院校 与社区学院模式}：小班研讨、个性化项目驱动与低门槛衔接通道经验互鉴。
    \item \textbf{拉美、阿拉伯世界与非洲的计算机教育崛起}：依托开源课程、低成本硬件与远程数字化赋能跨越式发展。
\end{itemize}

\subsection{专业建设与面向未来的思考}
\begin{itemize}
    \item \textbf{CS + X 跨学科复合专业}：计算生物学、数字法学、计算金融等新兴交叉学位建设范式。
    \item \textbf{智能时代的教学资源与教材变革}：从静态纸质教材向可交互 Jupyter/Colab 在线工坊、生成式 AI 交互助教演进。
\end{itemize}
